\chapter{Complexity — DAG Circuit Sat Lang}
%%% generated by the blueprint pipeline from TCSlib.Complexity.CircuitComplexity.DAGCircuitSatLang
\section{Overview}

Declarations of the Lean module \texttt{TCSlib.Complexity.CircuitComplexity.DAGCircuitSatLang}.

\section{Declarations}

\begin{definition}[Bounded list decoder]
\lean{BoolCircuit.decodeList}
\label{BoolCircuit.decodeList}
\leanok
Given a decoder that reads a Boolean list and may return an element of a type together
with the unread remainder, define a recursive decoder with a natural-number fuel bound.
Zero fuel and an empty input both fail; a leading false succeeds with the empty list and
the remaining input, while a leading true decodes one element and recursively decodes
the remainder with one less unit of fuel, failing if either decoding step fails.
\end{definition}

\begin{definition}[Gate kind decoder]
\lean{BoolCircuit.GateKind.decode}
\label{BoolCircuit.GateKind.decode}
\leanok
\uses{BoolCircuit.GateKind}
A Boolean list beginning with $00$, $01$, or $10$ decodes respectively to an AND, OR, or
NOT gate, together with the list after those first two bits. Every other input fails.
\end{definition}

\begin{definition}[DAG gate decoder]
\lean{BoolCircuit.DAGGate.decode}
\label{BoolCircuit.DAGGate.decode}
\leanok
\uses{BoolCircuit.DAGGate, BoolCircuit.GateKind.decode, BoolCircuit.decodeList, BoolCircuit.decodeNat}
Decode a Boolean list first as a gate kind and then, from the remaining bits, as a list
of natural-number arguments, using the length of that remaining list as the recursion
bound. If both decodings succeed, return the gate with that kind and argument list,
together with the final unread bits; otherwise fail.
\end{definition}

\begin{definition}[Validated circuit construction]
\lean{BoolCircuit.DAGCircuit.ofParts}
\label{BoolCircuit.DAGCircuit.ofParts}
\leanok
\uses{BoolCircuit.DAGCircuit, BoolCircuit.DAGGate}
Given a natural number $n$, a list of gates, and an output index, return the circuit
with those gates and output if every argument read by gate $i$ is less than $n+i$ and
the output index is less than $n$ plus the number of gates. If either condition fails,
return no circuit.
\end{definition}

\begin{definition}[Decoding a directed acyclic circuit]
\lean{BoolCircuit.DAGCircuit.decode}
\label{BoolCircuit.DAGCircuit.decode}
\leanok
\uses{BoolCircuit.DAGCircuit, BoolCircuit.DAGCircuit.encode, BoolCircuit.DAGCircuit.ofParts, BoolCircuit.DAGGate.decode, BoolCircuit.decodeList, BoolCircuit.decodeNat}
For a Boolean list $w$, the decoder returns a pair consisting of a natural number $n$
and a directed acyclic circuit on $n$ input vertices exactly when $w$ parses
successively as $n$, a list of gates, and an output index, the gates read only earlier
vertices, the output is a vertex of the circuit, and encoding the resulting circuit
reproduces $w$. In all other cases it returns no value.
\end{definition}

\begin{theorem}[List encoding and decoding]
\lean{BoolCircuit.decodeList_encodeList}
\label{BoolCircuit.decodeList_encodeList}
\leanok
\uses{BoolCircuit.decodeList, BoolCircuit.encodeList}
\difficulty{4}
Let $\alpha$ be a type, let $f$ be a decoder from Boolean lists to either no result or
an element of $\alpha$ together with a remaining Boolean list, and let $g$ assign a
Boolean list to each element of $\alpha$. For a finite list $A$ in $\alpha$, suppose
that for every $a\in A$ and every Boolean-list suffix $r$, decoding the concatenation of
$g(a)$ with $r$ returns $(a,r)$. Then for every natural-number fuel bound strictly
greater than the length of $A$ and every suffix $r$, decoding the encoding of $A$
followed by $r$ returns $(A,r)$.
\end{theorem}

\begin{theorem}[Decoding an encoded gate kind]
\lean{BoolCircuit.GateKind.decode_encode}
\label{BoolCircuit.GateKind.decode_encode}
\leanok
\uses{BoolCircuit.GateKind, BoolCircuit.GateKind.decode, BoolCircuit.GateKind.encode}
\difficulty{4}
For every gate kind $k$ and every Boolean list $r$, decoding the encoding of $k$
followed by $r$ returns $k$ together with the remaining list $r$.
\end{theorem}

\begin{theorem}[Decoding an encoded gate]
\lean{BoolCircuit.DAGGate.decode_encode}
\label{BoolCircuit.DAGGate.decode_encode}
\leanok
\uses{BoolCircuit.DAGGate, BoolCircuit.DAGGate.decode, BoolCircuit.DAGGate.encode, BoolCircuit.GateKind.decode_encode, BoolCircuit.decodeList_encodeList, BoolCircuit.decodeNat, BoolCircuit.decodeNat_encodeNat, BoolCircuit.encodeNat, BoolCircuit.length_encodeList_ge}
\difficulty{3}
For every Boolean circuit gate $g$ and every Boolean list $r$, decoding the encoding of
$g$ followed by $r$ returns exactly the pair $(g,r)$.
\end{theorem}

\begin{theorem}[Reconstruction from Circuit Parts]
\lean{BoolCircuit.DAGCircuit.ofParts_self}
\label{BoolCircuit.DAGCircuit.ofParts_self}
\leanok
\uses{BoolCircuit.DAGCircuit, BoolCircuit.DAGCircuit.ofParts}
\difficulty{1}
Let $n$ be a nonnegative integer and let $C$ be a Boolean circuit with $n$ input
vertices, gates in topological order, and an output vertex. Each gate at position $i$
reads only vertices with index less than $n+i$, and the output index is less than $n$
plus the number of gates. Constructing a circuit from the number of inputs, the gate
list of $C$, and its output index returns exactly $C$.
\end{theorem}

\begin{theorem}[Decoding an encoded circuit]
\lean{BoolCircuit.DAGCircuit.decode_encode}
\label{BoolCircuit.DAGCircuit.decode_encode}
\leanok
\uses{BoolCircuit.DAGCircuit, BoolCircuit.DAGCircuit.decode, BoolCircuit.DAGCircuit.encode, BoolCircuit.DAGCircuit.ofParts_self, BoolCircuit.DAGGate.decode, BoolCircuit.DAGGate.decode_encode, BoolCircuit.DAGGate.encode, BoolCircuit.decodeList_encodeList, BoolCircuit.decodeNat_encodeNat, BoolCircuit.encodeNat, BoolCircuit.length_encodeList_ge}
\difficulty{4}
For every finite acyclic Boolean circuit $C$ with $n$ inputs, decoding its
Boolean-string encoding returns exactly the pair consisting of $n$ and $C$.
\end{theorem}

\begin{theorem}[Encoding after successful decoding]
\lean{BoolCircuit.DAGCircuit.encode_of_decode}
\label{BoolCircuit.DAGCircuit.encode_of_decode}
\leanok
\uses{BoolCircuit.DAGCircuit, BoolCircuit.DAGCircuit.decode, BoolCircuit.DAGCircuit.encode}
\difficulty{4}
Let $w$ be a finite Boolean string, and let $C$ be a directed acyclic Boolean circuit
with some number of input vertices, gates in topological order, and a valid output
vertex. If decoding $w$ succeeds and returns $C$, then encoding $C$ recovers $w$.
\end{theorem}

\begin{definition}[DAG circuit satisfiability language]
\lean{BoolCircuit.dagCktSatLang}
\label{BoolCircuit.dagCktSatLang}
\leanok
\uses{BoolCircuit.DAGCircuit.encode, BoolCircuit.dagCktSat}
The language consists of exactly those Boolean lists that are the encodings of circuits
in the specified set of satisfiable DAG circuits.
\end{definition}

\begin{theorem}[Decoding the Circuit Satisfiability Language]
\lean{BoolCircuit.mem_dagCktSatLang_iff_decode}
\label{BoolCircuit.mem_dagCktSatLang_iff_decode}
\leanok
\uses{BoolCircuit.DAGCircuit.decode, BoolCircuit.DAGCircuit.decode_encode, BoolCircuit.DAGCircuit.encode_of_decode, BoolCircuit.dagCktSat, BoolCircuit.dagCktSatLang}
\difficulty{4}
A bit string belongs to the language of encodings of fan-in-two satisfiable Boolean DAG
circuits if and only if it decodes to a circuit and the resulting variable-count and
circuit pair belongs to the set of fan-in-two satisfiable DAG circuits.
\end{theorem}

\begin{theorem}[Membership in the satisfiable circuit language]
\lean{BoolCircuit.encode_mem_dagCktSatLang_iff}
\label{BoolCircuit.encode_mem_dagCktSatLang_iff}
\leanok
\uses{BoolCircuit.DAGCircuit, BoolCircuit.DAGCircuit.decode_encode, BoolCircuit.DAGCircuit.encode, BoolCircuit.dagCktSat, BoolCircuit.dagCktSatLang, BoolCircuit.mem_dagCktSatLang_iff_decode}
\difficulty{3}
For a natural number $n$ and a directed acyclic circuit $C$ on $n$ inputs, the encoding
of $C$ belongs to the language of encodings of fan-in-two satisfiable circuits if and
only if $C$ is fan-in-two and has an input assignment on which its output is true.
\end{theorem}

\begin{definition}[Circuit to CNF map]
\lean{BoolCircuit.dagCktSatToCNF}
\label{BoolCircuit.dagCktSatToCNF}
\leanok
\uses{BoolCircuit.DAGCircuit.IsFaninTwo, BoolCircuit.DAGCircuit.decode, BoolCircuit.DAGCircuit.toCNF}
Given a Boolean list, decode it as a DAG circuit. If decoding succeeds and every gate
has fan-in two, return the circuit's CNF; if decoding fails or some gate does not have
fan-in two, return the CNF consisting of the single empty clause.
\end{definition}

\begin{definition}[Circuit to 3-SAT serialization]
\lean{BoolCircuit.dagCktSatToSAT3}
\label{BoolCircuit.dagCktSatToSAT3}
\leanok
\uses{BoolCircuit.dagCktSatToCNF, Std.Sat.CNF.serialize}
Given a Boolean list, first map it to a CNF using the circuit-to-CNF map, then serialize
that CNF as a Boolean list.
\end{definition}

\begin{theorem}[Encoding of an empty clause is not in 3SAT]
\lean{BoolCircuit.not_mem_SAT3_serialize_nil}
\label{BoolCircuit.not_mem_SAT3_serialize_nil}
\leanok
\uses{Complexity.SAT3, Std.Sat.CNF.decode_serialize, Std.Sat.CNF.serialize}
\difficulty{3}
The encoding of the CNF formula consisting of a single empty clause does not belong to
the language of bit strings whose decoded formulas have width at most three and are
satisfiable.
\end{theorem}

\begin{theorem}[DAG Circuit Satisfiability and 3-SAT]
\lean{BoolCircuit.mem_dagCktSatLang_iff_mem_SAT3}
\label{BoolCircuit.mem_dagCktSatLang_iff_mem_SAT3}
\leanok
\uses{BoolCircuit.DAGCircuit.decode, BoolCircuit.dagCktSatLang, BoolCircuit.dagCktSatToCNF, BoolCircuit.dagCktSatToSAT3, BoolCircuit.mem_dagCktSatLang_iff_decode, BoolCircuit.mem_dagCktSat_iff, BoolCircuit.mem_dagCktSat_iff_toCNF, BoolCircuit.not_mem_SAT3_serialize_nil, Complexity.SAT3, Std.Sat.CNF.Satisfiable, Std.Sat.CNF.WidthAtMost, Std.Sat.CNF.decode, Std.Sat.CNF.decode_serialize}
\difficulty{5}
A bit string belongs to the language of encodings of fan-in-two satisfiable Boolean DAG
circuits if and only if the bit string produced by converting its circuit-satisfiability
CNF encoding to a bit string belongs to the language of satisfiable CNF formulas of
width at most three.
\end{theorem}

\begin{definition}[Boolean negation circuit]
\lean{BoolCircuit.notCircuit}
\label{BoolCircuit.notCircuit}
\leanok
\uses{BoolCircuit.DAGCircuit}
The one-input Boolean circuit consists of a single NOT gate reading input vertex $0$,
with output vertex $1$.
\end{definition}

\begin{definition}[Contradiction circuit]
\lean{BoolCircuit.contraCircuit}
\label{BoolCircuit.contraCircuit}
\leanok
\uses{BoolCircuit.DAGCircuit}
The one-input Boolean circuit has a NOT gate reading input vertex $0$, followed by an
AND gate reading vertices $0$ and $1$; its output is vertex $2$.
\end{definition}
